\documentclass[12pt,a4paper]{ctexart}
\usepackage{amsmath,amssymb}
\usepackage[margin=2.5cm]{geometry}
\usepackage{booktabs}
\usepackage{array}
 
\title{\textbf{计算机控制系统设计与实践\\ 第二周作业}}
\author{}
\date{}
 
\begin{document}
\maketitle
 
\noindent
姓\quad 名：\underline{\ 刘峻豪\ }\hspace{4em} 学\quad 号：\underline{\ 3230105220\ }\\[1.2ex]
专\quad 业：\underline{\ 自动化\ }\hspace{4em} 指导教师：\underline{\ 喻洁\ }
 
\bigskip
\hrule
\bigskip
 
\section*{题目}
\addcontentsline{toc}{section}{题目}
在设计和开发计算机控制系统时，检测仪表的选型需要关注哪些要点和注意事项？请简明扼要进行说明。
 
\section{概述}
检测仪表是计算机控制系统的“感官”，负责把温度、压力、流量、液位、成分等被测量转换为
可供计算机采集的电信号。其性能直接决定测量数据的\textbf{真实性、及时性与可靠性}，进而影响
控制品质与系统安全。因此选型时应结合\textbf{被测对象、工艺条件、系统接口、可靠性与经济性}
进行综合权衡，做到“够用、可靠、易维护、可集成”。
 
\section{选型要点}
\begin{enumerate}
  \item \textbf{被测参数与测量原理的匹配}：首先明确被测物理量及其变化范围，选择原理适用的
        传感器/变送器（如温度可选热电偶、热电阻、红外；流量可选电磁、涡街、差压等），
        使测量原理与介质特性、工况相适应。
  \item \textbf{量程与测量范围}：量程应覆盖被测量的正常波动及可能的极值，并留有适当裕度，
        使正常工作点落在量程的合理区段（一般约 $30\%\sim70\%$），既保证测量精度又防止超量程。
  \item \textbf{精度（准确度）等级}：按控制要求选择合适的精度等级，兼顾系统总误差分配；
        并非越高越好，过高精度会显著增加成本且未必带来控制效果的提升。
  \item \textbf{分辨率与灵敏度}：分辨率应与后续 A/D 转换位数、显示与控制精度相匹配，
        能分辨工艺所需的最小变化量。
  \item \textbf{动态特性（响应速度、时间常数）}：仪表的响应时间应快于被测量的变化和控制周期，
        避免因测量滞后引起控制品质下降或采样混叠。
  \item \textbf{输出信号形式与系统接口}：优先选用\textbf{标准化信号}（如 $4\sim20\,\mathrm{mA}$、
        $0\sim10\,\mathrm{V}$、HART、RS-485/Modbus、现场总线等），便于与 A/D、PLC/DCS 直接连接；
        远距离传输宜用电流信号或数字信号以抗干扰。
  \item \textbf{环境适应性与防护}：考虑现场温度、湿度、振动、粉尘、腐蚀气氛及电磁环境，
        选择相应的工作温度范围、防护等级（IP）、防爆等级与电磁兼容（EMC）性能。
  \item \textbf{被测介质特性}：针对介质的腐蚀性、粘度、脏污、导电性、温度与压力等，
        选择合适的材质、膜片、接液部件与安装方式，防止堵塞、结垢与损坏。
  \item \textbf{可靠性、稳定性与重复性}：关注平均无故障时间（MTBF）、长期零点/量程漂移、
        重复性与迟滞，保证长期运行数据可信。
  \item \textbf{供电与功耗}：核对供电方式（两线制/四线制、供电电压）与功耗，
        与系统电源和安全栅等相容。
  \item \textbf{安装、维护与校准便利性}：便于安装、拆卸、在线校准与故障排查，降低维护成本。
  \item \textbf{经济性与供货保障}：在满足要求的前提下兼顾性价比、备品备件、交货期与售后服务。
\end{enumerate}
 
\begin{table}[htbp]
  \centering
  \caption{检测仪表选型要点一览}
  \renewcommand{\arraystretch}{1.25}
  \begin{tabular}{>{\raggedright\arraybackslash}p{3.2cm} >{\raggedright\arraybackslash}p{9.2cm}}
    \toprule
    \textbf{要点} & \textbf{关注内容} \\
    \midrule
    测量原理     & 与被测量、介质、工况相适应 \\
    量程范围     & 覆盖工艺范围并留裕度，工作点居于量程合理区段 \\
    精度等级     & 满足控制要求、合理分配系统误差，避免过度追高 \\
    分辨率/灵敏度 & 与 A/D 位数、控制精度匹配 \\
    动态特性     & 响应快于被测量变化与控制周期 \\
    输出与接口   & 标准信号（$4\sim20\,\mathrm{mA}$、总线等），便于集成 \\
    环境/防护    & 温湿度、防爆、IP 等级、EMC 抗干扰 \\
    介质特性     & 腐蚀、粘度、脏污、温压对应的材质与结构 \\
    可靠性/稳定性 & MTBF、漂移、重复性、迟滞 \\
    供电/功耗    & 两/四线制、供电电压、与安全栅相容 \\
    安装/维护    & 安装位置代表性、在线校准、易维护 \\
    经济性       & 性价比、备件、供货与售后 \\
    \bottomrule
  \end{tabular}
\end{table}
 
\section{注意事项}
\begin{enumerate}
  \item \textbf{量程留裕度}：避免长期在量程上限附近或过低区段运行；正常工作点宜处于量程
        $30\%\sim70\%$ 左右，兼顾精度与安全。
  \item \textbf{精度的合理选择}：以满足工艺与控制要求为准，综合考虑传感器、变送器、A/D 及
        运算环节的\textbf{系统总精度}，不盲目追求单表高精度。
  \item \textbf{与采样和 A/D 的匹配}：仪表分辨率与量程应与 A/D 转换位数相匹配，减小量化误差；
        动态响应须满足采样定理，防止频率混叠。
  \item \textbf{信号传输与抗干扰}：远传优先电流或数字信号；采用屏蔽双绞线、单点接地、
        信号隔离等措施，抑制共模与电磁干扰，必要时加安全栅实现电气隔离。
  \item \textbf{安全与防爆}：易燃易爆或危险场所须选用本安/隔爆型并具备相应认证；
        关键回路考虑冗余与故障安全（fail-safe）设计。
  \item \textbf{安装位置的代表性}：测点应能真实反映被测量，避开死区、涡流、振动、强热源与
        管路弯头附近，保证测量代表性与稳定性。
  \item \textbf{校准与量值溯源}：投运前及周期性校准，保证量值可溯源，定期核查零点与量程漂移。
  \item \textbf{标准化与兼容}：统一信号制式与通信协议，便于与计算机控制系统集成、扩展与备件互换。
  \item \textbf{全生命周期成本}：综合考虑采购、安装、能耗、维护与更换等费用，而非仅看购置价。
\end{enumerate}
 
\section{小结}
检测仪表选型的核心是“\textbf{测得准、传得可靠、接得进、用得久}”：在准确匹配测量原理与工况的
基础上，合理确定量程与精度，保证动态特性与抗干扰能力，采用标准化输出便于与计算机系统集成，
并统筹可靠性、安全性与经济性，从而为计算机控制系统提供真实、及时、可信的测量信息。
 
\end{document}
